% =========================================================
% Proof: Co-Cauchy is an Equivalence Relation
% Source: notes/completeness/notes-co-cauchy.tex
% Difficulty: Medium
% =========================================================

\subsection*{Co-Cauchy is an Equivalence Relation}
\label{prf:co-cauchy-equivalence-relation}

\begin{remark*}[Return]
$\leftarrow$ \hyperref[prop:co-cauchy-equivalence-relation]{Back to Proposition in Notes}
\end{remark*}

\begin{proof}
\Claim The co-Cauchy relation $\sim$ on Cauchy sequences in $(X,d)$,
defined by $(x_n) \sim (y_n)$ iff $d(x_n, y_n) \to 0$,
is an equivalence relation.

\Given The co-Cauchy relation $\sim$ on the set of Cauchy sequences in $(X,d)$.

\Goal To verify reflexivity, symmetry, and transitivity.

\Strategy
\begin{itemize}
  \item \textbf{Reflexivity:} $d(x_n, x_n) = 0 \to 0$ by M2.
  \item \textbf{Symmetry:} $d(x_n, y_n) \to 0$ implies $d(y_n, x_n) \to 0$ by M3.
  \item \textbf{Transitivity:} If $d(x_n, y_n) \to 0$ and $d(y_n, z_n) \to 0$,
        bound $d(x_n, z_n) \le d(x_n, y_n) + d(y_n, z_n) \to 0$ by M4.
\end{itemize}

% -------------------------------------------------------
% YOUR PROOF HERE
% -------------------------------------------------------

\end{proof}

\begin{remark*}[Proof shape]
Three-part verification. Each part is one or two lines using a metric axiom.
\end{remark*}

\begin{remark*}[Dependencies]
\begin{itemize}
  \item Definition~\ref{def:metric}: M2 (identity), M3 (symmetry), M4 (triangle inequality).
  \item Definition of equivalence relation: reflexivity, symmetry, transitivity.
  \item Algebra of limits: sum of two sequences tending to $0$ tends to $0$.
\end{itemize}
\end{remark*}
